% !TeX spellcheck = <none>
\documentclass[12pt,a4paper]{article}

\usepackage[utf8]{inputenc}
\usepackage[T1]{fontenc}
\usepackage[french]{babel}

\usepackage{amsmath, amssymb}

\usepackage{graphicx}
\usepackage{caption}

\usepackage{siunitx} % Pour les unités

\usepackage{geometry}
\geometry{top=2cm, bottom=2.5cm, left=2cm, right=2cm}
\setlength{\parindent}{0pt}

\usepackage[sfdefault]{roboto}

\usepackage{enumitem}

\usepackage{hyperref}

\usepackage{xcolor}

\usepackage{multicol}
\setlength{\columnseprule}{1pt}
\def\columnseprulecolor{\color{gray}}


\usepackage{import}
\usepackage{xifthen}
\usepackage{pdfpages}
\usepackage{transparent}
\usepackage{float}

\newcommand{\incfig}[1]{%
	\def\svgwidth{\columnwidth}
	\import{./figures/}{#1.pdf_tex}
}



\title{Cours génie parasisique}
\author{}
\date{\today}

\begin{document}
	
	\maketitle
	
	\tableofcontents
	
	\newpage
	
	\part{Chapitre 1}
	
	\section{Partie 1 du cour}
	
	\subsection{Eléments du sismologie}
	
	La science qui étudie les séismes
	
	\subsubsection{Manifestaion des séismes}
	
	
	\begin{figure}[H]
		\centering
		\includegraphics[width=1.1\linewidth]{"figures/chapitre 1/Manifestaion des seismes/manifestaion des seismes"}
		\caption{Manifestaion des séismes}
		\label{fig:manifestaion-des-seismes}
	\end{figure}
	
	
	
	
	
	
	
	
	
	
	
	
	
	
	
	
	
	
	
	
	
	
	
	
	
	
	
	
	
	
	
	
	
\end{document}